\section{Formal verification}\label{sec:lean-proofs}

The main text contains the human-readable proofs of the barycentre
representation, the aggregate transfer inequality, the information-certified
cap, the sharp certificate, envelope exactness, certificate convexity, and the
pairwise bridge.
Where the bound is sharp, attainment of the optimal multi-marginal plan is
supplied as a premise; the Lean record does not formalize the general existence
theorem.

The formal record checks the displayed identities and bounds under their
stated assumptions.
These results concern the algebraic structure of the coherent portfolio
certificate and do not validate the article-embedding proxy, the economic
content of the carrier-and-slack restrictions, the empirical distance
estimator, or the chronological evaluation design.

The exact machine-checked scope and assumptions are:
\ppLeanStatusList

\section{Barycentre representation of dispersion}\label{sec:barycentre-appendix}

\begin{proposition}[Barycentre representation]
  \label{prop:p3-barycenter-representation}
  For laws $P_1,\dots,P_n$ on $\mathcal H$ with finite second moments,
  \eqref{eq:p3-barycenter-representation} holds.
\end{proposition}

For weights $q$ on the simplex, set
$\bar x_q:=\sum_i q_i x_i$.
The pointwise weighted-variance identity
\[
  \sum_{i<j}q_i q_j\lVert x_i-x_j\rVert^2
  =\sum_i q_i\lVert x_i-\bar x_q\rVert^2
\]
is the bridge between the two representations.
Given any multimarginal coupling of $P_1,\dots,P_n$, the random barycentre
$\bar X_q=\sum_i q_i X_i$ induces a law $Q$.
Applying the identity and then minimizing over multimarginal couplings gives
the direction from the joint-coupling infimum to the barycentre infimum.

Conversely, take couplings of $(P_i,Q)$ whose costs approach
$W_2^2(P_i,Q)$ and glue their conditional laws given the common
$Q$-distributed variable.
The resulting multimarginal coupling has the required marginals.
Conditionally on that common variable, the pointwise weighted-variance
minimum gives
\[
  \sum_{i<j}q_i q_j\lVert X_i-X_j\rVert^2
  \le \sum_i q_i\lVert X_i-Q\rVert^2.
\]
Taking expectations, then infima over the couplings and over $Q$, gives the
reverse inequality.
Thus the infimum values agree without assuming that either infimum is
attained.

\section{Data and W2 artifact construction}\label{sec:data-appendix}

The empirical universe contains \ppnum[0]{n-tickers} firms with both article
embeddings and compatible local return histories. Return alignment leaves
\ppnum[0]{n-obs} common daily observations.
Each article carries a creation date, URL hash, firm identifier, and a
4,096-dimensional Qwen3-Embedding-8B embedding. The representation normalizes
every row to unit Euclidean norm, and the ground cost is Euclidean chord
distance.

The registry distinguishes transport order from ground-cost exponent. The W2
entries use balanced quadratic transport, rooted normalization, and ground
exponent one. The persisted artifact is rooted W2; Paper 3 squares it once at
the certificate boundary, with regression tests enforcing that transition.

Balanced assignment currently requires equal empirical cloud sizes. Rows are
first restricted by article date and then ordered by URL hash. At each annual
cutoff, the governed artifact registry fixes a common cloud size supported by
every firm. This rule retains all \ppnum[0]{n-tickers} priced firms rather than
deleting thin-coverage firms; the early vintages consequently use fewer
articles per firm and should be read as noisier geometries.

The canonical priced W$_2$ roster used by the certificate and its diagnostics
is listed in \Cref{tab:sample-universe}.

% AUTO-GENERATED by pipeline render-sample-universe-table -- do not edit by hand.

\begin{longtable}{lll}
\caption{Canonical empirical roster: ticker symbol, firm name, and sector label from Nasdaq summary metadata for the firms in the canonical priced $W_2$ geometry, sorted by sector then symbol. Authors' calculations.}
\label{tab:sample-universe} \\
\toprule
Symbol & Name & Sector \\
\midrule
\endfirsthead
\toprule
Symbol & Name & Sector \\
\midrule
\endhead
\midrule
\endfoot
\bottomrule
\endlastfoot
BIDU & Baidu, Inc. & Communication Services \\
CHTR & Charter Communications, Inc. & Communication Services \\
CMCSA & Comcast Corporation & Communication Services \\
EA & Electronic Arts Inc. & Communication Services \\
GOOG & Alphabet Inc. & Communication Services \\
MTCH & Match Group Inc. & Communication Services \\
NFLX & Netflix Inc. & Communication Services \\
NTES & NetEase, Inc. & Communication Services \\
SIRI & Sirius XM Holdings Inc. & Communication Services \\
TMUS & T-Mobile US Inc. & Communication Services \\
TTWO & Take-Two Interactive Software, Inc. & Communication Services \\
BKNG & Booking Holdings Inc. & Consumer Cyclical \\
EBAY & eBay Inc. & Consumer Cyclical \\
EXPE & Expedia Group Inc. & Consumer Cyclical \\
HAS & Hasbro Inc. & Consumer Cyclical \\
JD & JD.com Inc. & Consumer Cyclical \\
LULU & Lululemon Athletica Inc. & Consumer Cyclical \\
MAR & Marriott International Inc. & Consumer Cyclical \\
MAT & Mattel Inc. & Consumer Cyclical \\
MELI & MercadoLibre Inc. & Consumer Cyclical \\
NCLH & Norwegian Cruise Line Holdings Ltd. & Consumer Cyclical \\
ORLY & O'Reilly Automotive Inc. & Consumer Cyclical \\
ROST & Ross Stores, Inc. & Consumer Cyclical \\
SBUX & Starbucks Corporation & Consumer Cyclical \\
TCOM & Trip.com Group Limited & Consumer Cyclical \\
TRIP & Tripadvisor, Inc. & Consumer Cyclical \\
TSCO & Tractor Supply Company & Consumer Cyclical \\
ULTA & Ulta Beauty Inc. & Consumer Cyclical \\
WYNN & Wynn Resorts Limited & Consumer Cyclical \\
COST & Costco Wholesale Corporation & Consumer Defensive \\
DLTR & Dollar Tree Inc. & Consumer Defensive \\
KDP & Keurig Dr Pepper Inc. & Consumer Defensive \\
KHC & The Kraft Heinz Company & Consumer Defensive \\
MDLZ & Mondelez International, Inc. & Consumer Defensive \\
MNST & Monster Beverage Corporation & Consumer Defensive \\
PEP & PepsiCo Inc. & Consumer Defensive \\
FANG & Diamondback Energy, Inc. & Energy \\
PYPL & PayPal Holdings Inc. & Financial Services \\
ALGN & Align Technology Inc. & Healthcare \\
AMGN & Amgen Inc. & Healthcare \\
AZN & AstraZeneca PLC & Healthcare \\
BIIB & Biogen Inc. & Healthcare \\
BMRN & BioMarin Pharmaceutical Inc. & Healthcare \\
DXCM & DexCom, Inc. & Healthcare \\
GILD & Gilead Sciences Inc. & Healthcare \\
HOLX & Hologic, Inc. & Healthcare \\
IDXX & IDEXX Laboratories, Inc. & Healthcare \\
ILMN & Illumina, Inc. & Healthcare \\
INCY & Incyte Corporation & Healthcare \\
ISRG & Intuitive Surgical Inc. & Healthcare \\
REGN & Regeneron Pharmaceuticals Inc. & Healthcare \\
VRTX & Vertex Pharmaceuticals Incorporated & Healthcare \\
AAL & American Airlines Group Inc. & Industrials \\
CSX & CSX Corporation & Industrials \\
FAST & Fastenal Company & Industrials \\
HON & Honeywell International Inc. & Industrials \\
JBHT & J.B. Hunt Transport Services, Inc. & Industrials \\
PCAR & PACCAR Inc & Industrials \\
UAL & United Airlines Holdings Inc. & Industrials \\
ADBE & Adobe Inc. & Technology \\
ADI & Analog Devices, Inc. & Technology \\
ADP & Automatic Data Processing Inc. & Technology \\
ADSK & Autodesk, Inc. & Technology \\
AKAM & Akamai Technologies, Inc. & Technology \\
AMAT & Applied Materials Inc. & Technology \\
AMD & Advanced Micro Devices Inc. & Technology \\
ASML & ASML Holding N.V. & Technology \\
AVGO & Broadcom Inc. & Technology \\
CDNS & Cadence Design Systems, Inc. & Technology \\
CHKP & Check Point Software Technologies Ltd. & Technology \\
CSCO & Cisco Systems Inc. & Technology \\
CTSH & Cognizant Technology Solutions Corporation & Technology \\
ENPH & Enphase Energy Inc. & Technology \\
FISV & Fiserv, Inc. & Technology \\
FTNT & Fortinet Inc. & Technology \\
INTC & Intel Corporation & Technology \\
INTU & Intuit Inc. & Technology \\
KLAC & KLA Corporation & Technology \\
LRCX & Lam Research Corporation & Technology \\
MCHP & Microchip Technology Incorporated & Technology \\
MRVL & Marvell Technology Inc. & Technology \\
MU & Micron Technology Inc. & Technology \\
NTAP & NetApp, Inc. & Technology \\
NVDA & NVIDIA Corporation & Technology \\
NXPI & NXP Semiconductors N.V. & Technology \\
OKTA & Okta, Inc. & Technology \\
PANW & Palo Alto Networks Inc. & Technology \\
PAYX & Paychex, Inc. & Technology \\
QCOM & QUALCOMM Incorporated & Technology \\
SNPS & Synopsys, Inc. & Technology \\
STX & Seagate Technology Holdings plc & Technology \\
SWKS & Skyworks Solutions, Inc. & Technology \\
TEAM & Atlassian Corporation & Technology \\
TXN & Texas Instruments Incorporated & Technology \\
WDAY & Workday Inc. & Technology \\
WDC & Western Digital Corporation & Technology \\
ZS & Zscaler Inc. & Technology \\
AEP & American Electric Power Company, Inc. & Utilities \\
EXC & Exelon Corporation & Utilities \\
XEL & Xcel Energy Inc. & Utilities \\
\end{longtable}


% Temporarily omitted while the chronological-validation
% implementation is audited.
\iffalse
\section{Chronological validation protocol}\label{sec:validation-appendix}

Formation occurs monthly. Marginal volatility uses 252 strictly prior trading
days, and realized variance uses the next 63 trading days. Calibration ends in
2022; the evaluation panel runs from 2023-01-04 through 2026-07-15, and its
feasible monthly formations run through 2026-04-01 so every one retains a full
63-day outcome. All use the frozen 2022 information geometry. Holdout returns
never enter transmission-cell selection.

Article uncertainty is evaluated with per-firm stationary time blocks. The
reported lower certificate endpoint is conditional on that resampling design
and on the fixed portfolio; it is not a portfolio-uniform confidence statement.
Return-block intervals govern calibration margin and coverage. The primary
comparison is inverse volatility, with equal weight, sample GMV, and
Ledoit--Wolf GMV as secondary benchmarks. Required residual slack and the
predeclared budgets $\delta\in\{0,0.025,0.05\}$ are sensitivity diagnostics.

\section{Full-sample numerical non-vacuity}\label{sec:nonvacuity-appendix}

The full-sample surface is retained only to show that the algebraic certificate
can be numerically nonzero before chronological validity is imposed. It uses
equal-weight and inverse-volatility portfolios fixed on the complete
2018--2022 panel and varies the prespecified transmission grid. It therefore
does not calibrate $(L,\tau)$, establish historical coverage, or explain
realized risk.

\Cref{fig:certificate-surface} plots $100\mathcal C(q)$ against common slack.
Increasing either $L$ or $\tau$ weakens the pairwise floor and lowers credit.
The floor-activation rate is the share of ordered firm pairs for which
$[W_{2,ij}/L-2\tau]_+$ is positive; positive-floor quantiles summarize its
squared magnitude conditional on activation. These quantities diagnose why a
surface becomes vacuous, but they are not sampling probabilities.

\begin{figure}[H]
  \centering
  \ppfigure{0.95}{certificate_surface}
  \caption{Full-sample certificate sensitivity. Each line reports
    $100\mathcal C(q)$ as common slack $\tau$ increases for one $L$. The
    calculation is descriptive numerical non-vacuity only and is not the
  frozen chronological validation.}
  \label{fig:certificate-surface}
\end{figure}

At $L=1$ and $\tau=0$, the deductions are
\ppnum[1]{certificate-equal-weight-least-conservative-pct}\% for equal weight
and \ppnum[1]{certificate-inverse-volatility-least-conservative-pct}\% for
inverse volatility. At $L=1.5$ and $\tau=0.10$, they are
\ppnum[1]{certificate-equal-weight-most-conservative-pct}\% and
\ppnum[1]{certificate-inverse-volatility-most-conservative-pct}\%.
Their numerical size cannot overturn the main result that no grid cell passed
the frozen calibration coverage rule.
\fi

\section{Supplementary empirical results}
\label{sec:p3-supplementary-results}

The main text reports the conventional benchmark values in
\Cref{tab:news-variance-benchmark}.
\Cref{fig:news-variance-benchmark} provides the corresponding visual comparison.

\begin{figure}[H]
  \centering
  \ppfigure{0.72}{news_variance_benchmark}
  \caption{Standardized in-sample variance relative to the full-sample
    long-only GMV benchmark. The news-only allocation is a descriptive
  zero-slack maximizer of the news certificate.}
  \label{fig:news-variance-benchmark}
\end{figure}

\subsection{Complete embedding-representation sensitivity}
\label{sec:p3-representation-ablation}

The sensitivity tables hold the priced firms, standardized-return covariance,
reference laws, caps, and bootstrap scheme fixed while changing the
representation.
For representation $r$, let
$I_r=100\,V(q_r)/V(q_{\mathrm{GMV}})$ denote its relative-GMV index.
Each paired estimand is $\Delta_{c,r}=I_c-I_r$ in index points, so a negative
contrast favours the candidate allocation on conventional variance.
The common return bootstrap re-standardizes returns and recomputes the GMV
denominator in each draw.
Consequently, a dash means that the news-only allocation violates the fixed
law's cap; the cap is not relaxed to manufacture a percentile.

% AUTO-GENERATED -- do not edit by hand.
\begin{table}[H]
\centering
\scriptsize
\setlength{\tabcolsep}{3.2pt}
\begin{tabular}{lrrrrrr}
\toprule
Encoder & Dim. & Uniform 12.5\% & Uniform 15\% & Fixed $\alpha=0.8$ & Concentrated 15\% & Rel. GMV \\
\midrule
Qwen3-8B full (primary) & 4096 & 0.10 & 0.08 & 0.19 & 0.46 & 147.5 \\
Qwen3-4B full & 2560 & 0.07 & 0.06 & 0.08 & 0.33 & 146.7 \\
Qwen3-8B@1024 & 1024 & 0.07 & 0.06 & 0.08 & 0.33 & 146.7 \\
Qwen3-4B@1024 & 1024 & 0.17 & 0.17 & 0.29 & 0.73 & 148.5 \\
BGE-large-v1.5 full & 1024 & -- & 0.01 & 0.03 & 0.18 & 145.3 \\
Qwen3-8B@256 & 256 & 0.27 & 0.29 & 0.46 & 0.98 & 149.4 \\
Qwen3-8B@64 & 64 & 0.01 & 0.01 & 0.01 & 0.13 & 144.6 \\
EttaX V0 & 320 & -- & 0.33 & 0.53 & 1.10 & 149.7 \\
EttaX V1 & 320 & 1.19 & 1.10 & 1.59 & 2.60 & 152.0 \\
EttaX V3 & 320 & -- & 0.77 & 1.20 & 2.05 & 151.4 \\
\midrule
Equal risk weights & -- & 34.73 & 34.02 & 32.84 & 28.33 & 162.6 \\
\bottomrule
\end{tabular}
\caption{Representation sensitivity of the news-only allocation across seven base representations and three matched EttaX encoder vintages on the common 100-firm, 1,207-date standardized-return panel. The four middle columns report the percentage of fixed-law allocations with standardized variance no greater than the candidate's; lower is better. Relative variance indexes the long-only sample GMV to 100, so values above 100 measure percentage excess variance and lower is better. A dash marks a candidate that violates the law's unchanged cap. The $\alpha=0.8$ law is a fixed moderately concentrated comparison; its realized mean effective $N$ is reported rather than assumed to match the primary allocation. Returns evaluate every fixed news-only allocation in sample; they do not construct it.}
\label{tab:p3-representation-ablation}
\end{table}
\medskip
\begin{table}[H]
\centering
\scriptsize
\setlength{\tabcolsep}{3.0pt}
\begin{tabularx}{\linewidth}{Xrrrrrr}
\toprule
Contrast ($c-r$) & Estimate & 95\% CI & Raw $p$ & Holm $p$ & Bound & Equiv. \\
\midrule
Qwen3-4B full $-$ Qwen3-8B full & -0.842 & [-2.664, 1.649] & 0.454 & 0.873 & -- & -- \\
Qwen3-4B@1024 $-$ Qwen3-8B@1024 & 1.811 & [-0.863, 3.486] & 0.192 & 0.768 & -- & -- \\
Qwen3-8B@1024 $-$ Qwen3-8B full & -0.794 & [-1.934, 0.762] & 0.291 & 0.873 & -- & -- \\
Qwen3-8B@256 $-$ Qwen3-8B@1024 & 2.641 & [0.076, 4.830] & 0.046 & 0.322 & -- & -- \\
Qwen3-8B@64 $-$ Qwen3-8B@256 & -4.774 & [-8.324, 0.365] & 0.065 & 0.390 & -- & -- \\
BGE-large-v1.5 full $-$ Qwen3-8B@1024 & -1.446 & [-3.807, 1.674] & 0.349 & 0.873 & -- & -- \\
BGE-large-v1.5 full $-$ Qwen3-4B@1024 & -3.257 & [-5.526, 0.733] & 0.115 & 0.575 & -- & -- \\
\addlinespace[1.5pt]
EttaX V1 $-$ EttaX V0 & 2.365 & [-2.930, 5.212] & 0.370 & 1.000 & 3.807 & No \\
EttaX V3 $-$ EttaX V0 & 1.693 & [-2.465, 4.096] & 0.467 & 1.000 & 3.807 & No \\
EttaX V3 $-$ EttaX V1 & -0.673 & [-2.713, 1.793] & 0.599 & 1.000 & 3.807 & Yes \\
\bottomrule
\end{tabularx}
\caption{Paired stationary-return-bootstrap contrasts for the ten-cell representation ladder. The seven base contrasts are descriptive candidate-minus-reference relative-GMV index-point differences, so a negative estimate favours the candidate. Raw $p$-values are two-sided add-one sign-tail probabilities; Holm adjustment is applied separately to the seven base and three vintage contrasts. The EttaX rows additionally use the BGE-large-minus-Qwen3-8B@1024 interval as a benchmark-calibrated, post-specified robustness threshold. Equivalence requires strict containment of the full interval; this is not a causal or architecture-only comparison.}
\label{tab:p3-representation-contrasts}
\end{table}


The full-width-to-1,024 Qwen3-8B contrast includes zero, so moderate truncation
does not reject zero difference in relative variance.
At fixed 1,024 width, the Qwen3-4B comparison and the 64-coordinate
Qwen3-8B stress-cell comparison with 256 coordinates do not survive Holm
adjustment.
Among the EttaX vintages, only the V3--V1 interval satisfies the stated
post-specified equivalence bound; neither comparison involving V0 does.
